\documentclass{article}
\usepackage[margin=1in]{geometry}
\usepackage{amsmath,amssymb}
\usepackage[most]{tcolorbox}
\usepackage[hidelinks]{hyperref}

\newcommand{\BookName}{Pattern Recognition and Machine Learning}
\newcommand{\BookAuthor}{Christopher M. Bishop}
\newcommand{\ChapterName}{Introduction}
\newcommand{\ExerciseID}{1.10}

\title{Chapter: \ChapterName}
\author{Ahonakpon Guy Gbaguidi}
\date{}

\begin{document}

\maketitle

\begin{center}
  \large\textbf{Exercise \ExerciseID}

  \vspace{0.5em}

  \normalsize\BookName\ - \BookAuthor
\end{center}

\begin{tcolorbox}[breakable,colback=gray!5,colframe=gray!60,title=Solution]
Let \(X\) and \(Z\) be two random variables. In this exercise, we want to show that, when
\(X\) and \(Z\) are statistically independent, the mean and variance of their sum satisfy
the following properties:

\begin{align*}
  \mathbb{E}[X+Z] &= \mathbb{E}[X] + \mathbb{E}[Z] \tag{1.128} \\
  \operatorname{var}[X+Z] &= \operatorname{var}[X] + \operatorname{var}[Z]. \tag{1.129}
\end{align*}

\subsubsection*{Proof of the expectation identity}

First, notice an important point: the identity
\[
  \mathbb{E}[X+Z] = \mathbb{E}[X] + \mathbb{E}[Z]
\]
does \emph{not} require \(X\) and \(Z\) to be independent. It only requires the relevant
expectations to exist. This property is called the linearity of expectation.

\paragraph{Continuous case.}
Assume that \(X\) and \(Z\) are continuous random variables with joint density
\(p(x,z)\). The joint density describes the probability of the pair \((X,Z)\). It may or
may not factorize as \(p(x)p(z)\). If \(X\) and \(Z\) are independent, then
\(p(x,z)=p(x)p(z)\), but we will not need that here.

By definition, the expectation of \(X+Z\) is
\[
  \mathbb{E}[X+Z]
  =
  \iint_{\mathbb{R}^2}
  (x+z)p(x,z)\,dx\,dz.
\]

Using the linearity of integration, we split the integrand:

\begin{align*}
  \mathbb{E}[X+Z]
  &=
  \iint_{\mathbb{R}^2}
  (x+z)p(x,z)\,dx\,dz \\
  &=
  \iint_{\mathbb{R}^2}
  xp(x,z)\,dx\,dz
  +
  \iint_{\mathbb{R}^2}
  zp(x,z)\,dx\,dz.
\end{align*}

Now we show that the first double integral is \(\mathbb{E}[X]\). For a joint density,
the marginal density of \(X\) is obtained by integrating out \(z\):
\[
  p_X(x) = \int_{-\infty}^{\infty} p(x,z)\,dz.
\]
Therefore,

\begin{align*}
  \iint_{\mathbb{R}^2}
  xp(x,z)\,dx\,dz
  &=
  \int_{-\infty}^{\infty}
  x
  \left(
    \int_{-\infty}^{\infty} p(x,z)\,dz
  \right)
  dx \\
  &=
  \int_{-\infty}^{\infty} x p_X(x)\,dx \\
  &=
  \mathbb{E}[X].
\end{align*}

Similarly, the marginal density of \(Z\) is
\[
  p_Z(z) = \int_{-\infty}^{\infty} p(x,z)\,dx,
\]
and so

\begin{align*}
  \iint_{\mathbb{R}^2}
  zp(x,z)\,dx\,dz
  &=
  \int_{-\infty}^{\infty}
  z
  \left(
    \int_{-\infty}^{\infty} p(x,z)\,dx
  \right)
  dz \\
  &=
  \int_{-\infty}^{\infty} z p_Z(z)\,dz \\
  &=
  \mathbb{E}[Z].
\end{align*}

Combining the two pieces gives
\[
  \mathbb{E}[X+Z]
  =
  \mathbb{E}[X]+\mathbb{E}[Z].
\]

The key point is that we never used \(p(x,z)=p_X(x)p_Z(z)\). Hence the result holds
whether \(X\) and \(Z\) are independent or dependent.

\paragraph{Discrete case.}
Now assume that \(X\) and \(Z\) are discrete random variables. Let
\[
  p(x,z) = \mathbb{P}(X=x,Z=z)
\]
be their joint probability mass function. Again, this joint probability may or may not
factorize into \(\mathbb{P}(X=x)\mathbb{P}(Z=z)\).

By definition,
\[
  \mathbb{E}[X+Z]
  =
  \sum_x \sum_z (x+z)\mathbb{P}(X=x,Z=z).
\]

Using the linearity of summation, we get

\begin{align*}
  \mathbb{E}[X+Z]
  &=
  \sum_x \sum_z x\mathbb{P}(X=x,Z=z)
  +
  \sum_x \sum_z z\mathbb{P}(X=x,Z=z).
\end{align*}

For the first term, \(x\) does not depend on \(z\), so

\begin{align*}
  \sum_x \sum_z x\mathbb{P}(X=x,Z=z)
  &=
  \sum_x x \sum_z \mathbb{P}(X=x,Z=z) \\
  &=
  \sum_x x \mathbb{P}(X=x) \\
  &=
  \mathbb{E}[X].
\end{align*}

The equality
\[
  \sum_z \mathbb{P}(X=x,Z=z) = \mathbb{P}(X=x)
\]
is simply the definition of the marginal probability of \(X\). We add up all possible
values of \(Z\), leaving only the probability that \(X=x\).

For the second term, \(z\) does not depend on \(x\), so

\begin{align*}
  \sum_x \sum_z z\mathbb{P}(X=x,Z=z)
  &=
  \sum_z z \sum_x \mathbb{P}(X=x,Z=z) \\
  &=
  \sum_z z \mathbb{P}(Z=z) \\
  &=
  \mathbb{E}[Z].
\end{align*}

Thus,
\[
  \mathbb{E}[X+Z]
  =
  \mathbb{E}[X]+\mathbb{E}[Z].
\]

Again, no independence assumption was used. The result follows only from the linearity
of sums and from the definition of marginal probabilities.

\subsubsection*{Proof of the variance identity}

For the variance identity, independence does matter. By definition,
\[
  \operatorname{var}[Y] = \mathbb{E}\left[(Y-\mathbb{E}[Y])^2\right].
\]
Taking \(Y=X+Z\), and using the expectation identity proved above, we have

\begin{align*}
  \operatorname{var}[X+Z]
  &=
  \mathbb{E}\left[\left(X+Z-\mathbb{E}[X+Z]\right)^2\right] \\
  &=
  \mathbb{E}\left[\left(X+Z-\mathbb{E}[X]-\mathbb{E}[Z]\right)^2\right] \\
  &=
  \mathbb{E}\left[\left((X-\mathbb{E}[X])+(Z-\mathbb{E}[Z])\right)^2\right].
\end{align*}

Now expand the square:

\begin{align*}
  \operatorname{var}[X+Z]
  &=
  \mathbb{E}\left[(X-\mathbb{E}[X])^2\right]
  +
  \mathbb{E}\left[(Z-\mathbb{E}[Z])^2\right] \\
  &\qquad
  +
  2\mathbb{E}\left[(X-\mathbb{E}[X])(Z-\mathbb{E}[Z])\right] \\
  &=
  \operatorname{var}[X]
  +
  \operatorname{var}[Z]
  +
  2\operatorname{cov}[X,Z].
\end{align*}

If \(X\) and \(Z\) are independent, then
\[
  \operatorname{cov}[X,Z]=0.
\]
Therefore,
\[
  \operatorname{var}[X+Z]
  =
  \operatorname{var}[X]+\operatorname{var}[Z].
\]

This completes the proof of both identities. The expectation identity holds for
independent and dependent random variables, while the variance identity needs
independence, or more generally zero covariance.
\end{tcolorbox}

\end{document}
